\section*{Strong Seeded Dispersers and Friends (original Sections 5.1 and 5.3)}
\subsection*{5.1 Strong Seeded Dispersers}
\begin{definition}[strong disperser]
A function $\Disp:\{0,1\}^n\times\{0,1\}^d\to\{0,1\}^m$ is a strong $(k,\varepsilon)$ disperser, if for every $(n,k)$ source $X$,
\[|\Supp((Y,\Disp(X,Y)))|>(1-\varepsilon)DM.\]
We say $\Disp$ is (source) linear if for every $y\in\{0,1\}^d$ and every $x_1,x_2\in\mathbb F_2^n$, $\Disp(x_1+x_2,y)=\Disp(x_1,y)+\Disp(x_2,y)$.
\end{definition}
We are interested in the important special case where $m=1$. In this case, non-explicitly, a random function is (w.h.p.) a strong $(k,\varepsilon)$ disperser with $d=\log n+\log(\frac1\varepsilon)+O(1)$ provided that $k\ge\log\log(\frac1\varepsilon)+O(1)$ [37,33]. A matching lower bound, up to additive constant factors, was given by [37].
Using the translation between strong seeded dispersers and erasure list-decodable codes which we discuss in Section 5.3, Guruswami and Indyk's result [26] gives a probabilistic polynomial time algorithm that outputs with high probability a strong seeded disperser with seed-length $d=2\log(\frac1\varepsilon)+\log n+\log\log(\frac1\varepsilon)$ and optimal entropy-loss. The construction can be made deterministic, but with running time exponential in $1/\varepsilon$.
Note that as we discuss the one output bit case, the required entropy is essentially the entropy-loss. From Theorem 27 we can derive a better explicit construction of a strong disperser with small error.
\begin{theorem}
For every constant $0<\gamma<1$ there exists a constant $c\ge1$ such that for every integer $n$ and $\varepsilon\le n^{-\frac c{1-\gamma}}$ there exists an explicit strong $(k,\varepsilon)$ disperser $\Disp:\{0,1\}^n\times\{0,1\}^d\to\{0,1\}$ where $d=(1+\gamma)\log(\frac1\varepsilon)$ and $k=c\log d$.
\end{theorem}
\begin{proof}
Set $\varepsilon_0=\frac14$ and $\delta_0=\frac\gamma{1+\gamma}$. Let $c$ be the constant from Theorem 27 for $\delta_0$ and $\varepsilon_0$ and let $\twoExt:[N]\times[D]\to\{0,1\}$ be the $((n,k),(d,k_2=\delta_0d),\varepsilon_0)$ two-source extractor where $d=(1+\gamma)\log(\frac1\varepsilon)$ and $k=c\log d$. Notice that $d\ge c\log n$ (because $\varepsilon\le n^{-\frac c{1+\gamma}}$) as required. Let $\Disp(x,y)=\twoExt(x,y)$.
Let $X\subseteq[N]$ be a set of size $K$ and call a value $y\in[D]$ $b$-bad if $\Disp(X,y)=\{b\}$. It follows that the sets of 0-bad $y$-s and 1-bad $y$-s are each of size less than $K_2$. Therefore,
\begin{align*}
|\Supp((U_d,\Disp(X,U_d)))|&>2K_2+2(D-2K_2)=2D-2K_2\\
&=\left(1-\frac{K_2}D\right)2D=(1-\varepsilon)2D,
\end{align*}
because $\frac{K_2}D=2^{-(1-\delta_0)d}=2^{-\frac1{1+\gamma}d}=2^{-\log(\frac1\varepsilon)}=\varepsilon$.
\end{proof}
\subsection*{5.3 Erasure List-Decodable Codes}
An $(\bar n,n)$ (binary) code is a mapping $\mathcal C:\{0,1\}^n\to\{0,1\}^{\bar n}$. The code $\mathcal C$ is linear if $\mathcal C$ is linear, and is denoted by $[\bar n,n]$. We identify a code with the image of $\mathcal C$. For a linear $\mathcal C$ this image is a linear subspace of $\mathbb F_2^{\bar n}$ of dimension $n$. A generator matrix for an $[\bar n,n]$ code $\mathcal C$ is any matrix whose columns form a basis for $\mathcal C$. In the erasures noise model, an adversarially chosen subset of the codeword's symbols are erased and the positions where erasures have occurred are known.
\begin{definition}[erasure list-decodable code]
A code $\mathcal C\subseteq\{0,1\}^{\bar n}$ is $(s,L)$ erasure list-decodable if for every $r\in\{0,1\}^{\bar n-s}$ and every set $T\subseteq[\bar n]$ of size $\bar n-s$,
\[|\{c\in\mathcal C\mid c|_T=r\}|<L,\]
where $c|_T$ denotes the projection of $c$ to the coordinates in $T$.
\end{definition}
The following folklore lemma (see, e.g., [24, Lemma 1]) gives an alternative characterization of linear erasure list-decodable codes.
\begin{lemma}
An $[\bar n,n]_2$ linear code $\mathcal C$ is $((1-\varepsilon)\bar n,L)$ erasure list-decodable if and only if its $\bar n\times n$ generator matrix $G$ has the property that every $\varepsilon\bar n\times n$ sub-matrix of $G$ has rank greater than $n-\log L$.
\end{lemma}
Non-explicitly, we have:
\begin{theorem}[{[24]}]
For every $n$ and $\varepsilon>0$, there exists an $(\bar n,n)$ binary code that is $((1-\varepsilon)\bar n,L)$-erasure list-decodable of rate $\frac n{\bar n}=\Omega(\varepsilon)$ and $L=O(\log(1/\varepsilon))$.
\end{theorem}
Guruswami [25] observed that strong dispersers can be used to construct erasure list-decodable codes. Here we complement his argument, and note that strong dispersers are equivalent to erasure list-decodable codes. Given a function $\Disp:[N]\times[D]\to\{0,1\}$, we consider the $(D,n)$ code $\mathcal C_{\Disp}:\{0,1\}^n\to\{0,1\}^D$ defined by $\mathcal C_{\Disp}(x)_i=\Disp(x,i)$. Note that the code is linear if and only if $\Disp$ is linear.
\begin{lemma}[following {[25, Lemma 12]}]
The function $\Disp:[N]\times[D]\to\{0,1\}$ is a strong $(k,\varepsilon)$ disperser if and only if $\mathcal C_{\Disp}$ is $((1-2\varepsilon)D,K)$ erasure list-decodable.
\end{lemma}
\begin{proof}
For one direction, assume $\Disp$ is a strong $(k,\varepsilon)$ disperser. We wish to prove that $\mathcal C_{\Disp}$ is $((1-2\varepsilon)D,K)$ erasure list-decodable. Let $T=\{t_1,\ldots,t_{2\varepsilon D}\}\subseteq[D]$ be an arbitrary set of size $2\varepsilon D$ and $r\in\{0,1\}^{2\varepsilon D}$ an arbitrary string. Let $X_{T,r}\subseteq\{0,1\}^n$ denote the set of all the messages $x$ for which $\mathcal C_{\Disp}(x)|_T=r$. Then,
\[|\Supp((U_d,\Disp(X_{T,r},U_d)))|\le|T|\cdot1+(D-|T|)\cdot2\le(1-\varepsilon)2D,\]
where the first inequality follows by considering seeds in $T$ and seeds in $[D]\setminus T$. For a seed $t_i\in T$ we have that $\Disp(X_{T,r},t_i)$ is fixed, hence each such seed contributes 1 to the support size. For any other seed $y$, the support size of $\Disp(X_{T,r},y)$ is at most 2. As $\Disp$ is a strong $(k,\varepsilon)$ disperser, we conclude that $|X_{T,r}|\le K$ as desired.
For the other direction assume $\Disp$ is a not a strong $(k,\varepsilon)$ disperser. Then, there exists a set $X\subseteq\{0,1\}^n$ such that $|X|\ge K$ and $|\Supp((U_d,\Disp(X,U_d)))|\le(1-2\varepsilon)2D$. Note that for every $y\in[D]$ we have $|\Supp(\Disp(X,y))|\in\{1,2\}$. Therefore, following the above calculation, there exists a set $T\subseteq D$ of size at least $2\varepsilon D$ such that for each $y\in T$, $|\Supp(\Disp(X,y))|=1$. But this means that for every $x\in X$, $\mathcal C_{\Disp}(x)|_T$ is the same (punctured) codeword. It follows that $\mathcal C_{\Disp}$ is not $((1-2\varepsilon)D,K)$ erasure list-decodable.
\end{proof}
\begin{corollary}
If $\Disp:\{0,1\}^n\times\{0,1\}^d\to\{0,1\}$ is a strong $(k,\varepsilon)$ disperser with seed-length $d=a_1\log n+a_2\log(\frac1\varepsilon)+a_3$ (for some $a_1\ge1,a_2\ge1$ and $a_3$) then $\mathcal C_{\Disp}$ is a $((1-2\varepsilon)D,K)$ erasure list-decodable code of rate $2^{-a_3}\cdot n^{1-a_1}\cdot\varepsilon^{a_2}$.
\end{corollary}
When $\varepsilon$ is much smaller than $\frac1n$ the dominant factor is determined by $a_2$. As we mentioned earlier (and as Guruswami also notes in [25]) previous explicit constructions for binary codes had $a_2\ge2$ (usually inherited from extractor constructions). Our construction is the first to get arbitrary close to $a_2=1$ and small list-size. Combining Corollary 45 and Theorem 40, we obtain:
\begin{theorem}
For every constant $0<\gamma<1$ there exists a constant $c\ge1$ such that for every integer $n$ and $\varepsilon\le n^{-\frac c{1-\gamma}}$ there exists an explicit code $\mathcal C:\{0,1\}^n\to\{0,1\}^{(\frac1\varepsilon)^{1+\gamma}}$ that is
\[\left((1-2\varepsilon)\left(\frac1\varepsilon\right)^{1+\gamma},\left((1+\gamma)\log\frac1\varepsilon\right)^c\right)\]
erasure list-decodable of rate $n\varepsilon^{1+\gamma}$.
\end{theorem}
